% Generated draft: requires mathematical and source review.
\documentclass[11pt]{article}
\usepackage[margin=25mm]{geometry}
\usepackage{fontspec}
\setmainfont{Latin Modern Roman}
\usepackage{amsmath,amssymb,mathtools,unicode-math}
\setmathfont{Latin Modern Math}
\usepackage{enumitem,needspace,xurl,hyperref}
\setlength{\parindent}{0pt}
\setlength{\parskip}{5pt}
\setlength{\emergencystretch}{3em}
\newcommand{\N}{\mathbb{N}}
\newcommand{\Z}{\mathbb{Z}}
\newcommand{\Q}{\mathbb{Q}}
\newcommand{\R}{\mathbb{R}}
\newcommand{\C}{\mathbb{C}}
\newcommand{\F}{\mathbb{F}}
\newcommand{\A}{\mathbb{A}}
\newcommand{\PP}{\mathbb{P}}
\newcommand{\E}{\mathbb{E}}
\newcommand{\eps}{\varepsilon}
\newcommand{\dd}{\,\mathrm{d}}
\begin{document}
\section*{UnsolvedMath ID: 10400204}
\textbf{Source problem code:} catalogue entry
\subsection{Definitions and mathematical statement}
Let \(M\) be a \(3\)-manifold in the source's domain for the LMO invariant, and let \(\xi\in H^1(M;\Z)\) be an integral degree-one cohomology class. T. Le and V. Turaev's Problem 11.8 asks to construct an extension assigning an invariant
\[
 Z_{\rm LMO}(M,\xi)
 \qquad\text{for every }\xi\in H^1(M;\Z),
\]
with the recovery condition
\[
 Z_{\rm LMO}(M,0)=Z_{\rm LMO}(M).
\]
The invariant should be attached to the pair, with the usual compatibility under a manifold equivalence carrying one cohomology class to the other; the equivalence category and any varying target-space identifications are not specified. The target, normalization and additional properties of the extension are absent from the extract.

The source's motivation is that the usual LMO invariant corresponds only to the trivial cohomology class and vanishes for sufficiently high first Betti number \(b_1(M)=\dim_{\Q}H^1(M;\Q)\). It does not state the threshold or normalization under which that vanishing claim is meant. It further reports that K. Habiro has an extension that \emph{might} solve the problem, without giving its definition or a reference. These qualifications do not establish a solution.

The zero-class identity alone permits an assignment independent of \(\xi\), so a precise intended problem needs the missing criterion for an informative cohomology-dependent extension. No such nontriviality, universality, finite-type, or gluing condition is invented here.
\subsection{Short English statement}
Extend the LMO invariant to three-manifolds equipped with every integral first-cohomology class, recovering the original invariant at the zero class. The target and required dependence on the class are unspecified. The supplied source mentions a possibly relevant Habiro construction without claiming it resolves the request.
\subsection{Sources}
\begin{itemize}
\item Local UnsolvedMath registry, permanent ID \(10400204\), exact source code \texttt{catalogue entry}, Problem 11.8, attributed to T. Le and V. Turaev. The statement also names K. Habiro's possibly relevant extension. The title is truncated after ``(T.''
\item Ohtsuki, \emph{Problems on invariants of knots and 3-manifolds} (2002), Problem 11.8, PDF page 157: \url{https://msp.org/gtm/2002/04/gtm-2002-04-024s.pdf}. Bold headings distinguish 231 numbered items from surrounding references.
\item Extraction pdf-font-structure, accessed 2026-07-29. Source evidence NEEDS\_REVIEW presents an open problem and calls for a current primary-literature audit. Registry status open and research\_classification OPEN-TRIAGE coexist with a research summary labeled PARTIAL-PROGRESS. Both source and summary say Habiro's extension may address the question; no full citation or definitive identification is supplied.
\item Dataset attribution: Ulam AI / UnsolvedMath Contributors, CC BY 4.0. Source rights NEEDS\_REVIEW: publicly accessible author/publisher copy, redistribution terms requiring review.
\end{itemize}
\end{document}
